\VerdierPhysicalPageAnchor{0061}%
\VerdierSemanticAnchor{ch1-1-6-2}
\VerdierSemanticAnchor{ch1-1-6-2-proposition}
\noindent\textbf{Proposition 1.6.2.}\quad

\VerdierSemanticAnchor{ch1-1-6-2-part1}
\noindent\emph{1) Let
$\alpha\in Z^2(\mathbb Z^n,\mathbb Z/2\mathbb Z)$ be a normalized
inhomogeneous cocycle, let
$\dot\alpha\in H^2(\mathbb Z^n,\mathbb Z/2\mathbb Z)$ be its cohomology
class, and let
$\dot\alpha_{i,j}\in H^2(\mathbb Z^n_{i,j},\mathbb Z/2\mathbb Z)$ be
its components. Then (using (1.6.1.2)):}

\VerdierSemanticAnchor{ch1-1-6-2-eq1}
\begin{equation}
  \dot\alpha_{i,j}
  =\alpha(e_i,e_j)^{-1}\alpha(e_j,e_i).
  \tag{1.6.2.1}
\end{equation}

\VerdierSemanticAnchor{ch1-1-6-2-convention}
\noindent\emph{(The group law on $\mathbb Z/2\mathbb Z$ is written
multiplicatively.)}

\VerdierSemanticAnchor{ch1-1-6-2-part2}
\noindent\emph{2) Let $\alpha_1$ and $\alpha_2$ be two cohomologous
elements of $Z^2(\mathbb Z^n,\mathbb Z/2\mathbb Z)$ and, for every
$i\in[1,n]$, let $\beta_i$ be an element of $\mathbb Z/2\mathbb Z$.
There exists a unique map}

\VerdierSemanticAnchor{ch1-1-6-2-beta-map}
\begin{equation*}
  \beta\colon\mathbb Z^n\longrightarrow\mathbb Z/2\mathbb Z
\end{equation*}
\noindent\emph{such that:}

\VerdierSemanticAnchor{ch1-1-6-2-eq2}
\begin{equation}
  \begin{aligned}
    \alpha_2(\sigma_1,\sigma_2)\alpha_1^{-1}(\sigma_1,\sigma_2)
      &=\beta(\sigma_2)\beta(\sigma_1+\sigma_2)^{-1}\beta(\sigma_1),\\[-1mm]
    &\hspace{7em}\text{for }\sigma_1,\sigma_2\in\mathbb Z^n,
  \end{aligned}
  \tag{1.6.2.2}
\end{equation}

\VerdierSemanticAnchor{ch1-1-6-2-eq3}
\begin{equation}
  \beta(e_i)=\beta_i.
  \tag{1.6.2.3}
\end{equation}

\VerdierSemanticAnchor{ch1-1-6-2-proof}
\VerdierSemanticAnchor{ch1-1-6-2-proof-first}
To prove the first assertion, we reduce to the case $n=2$ and
immediately verify that the quantity
$\alpha(e_i,e_j)^{-1}\alpha(e_j,e_i)$ does not change when $\alpha$ is
modified by a coboundary. When $\alpha$ is a coboundary, it is
symmetric, and consequently formula (1.6.2.1) holds in this case. On
the other hand, the cocycle

\VerdierSemanticAnchor{ch1-1-6-2-alpha0}
\begin{equation*}
  \alpha_0(p_1e_i+q_1e_j,p_2e_i+q_2e_j)=(-1)^{q_1p_2}
\end{equation*}
is not cohomologous to $0$, since

\VerdierSemanticAnchor{ch1-1-6-2-alpha0-test}
\begin{equation*}
  \alpha_0(e_i,e_j)^{-1}\alpha_0(e_j,e_i)=-1,
\end{equation*}
and every cocycle not cohomologous to zero is cohomologous to
$\alpha_0$.

\VerdierSemanticAnchor{ch1-1-6-2-proof-second}
To prove the second assertion, observe that two solutions of (1.6.2.2)
differ by a group homomorphism; consequently, one can always modify a
solution of (1.6.2.2), uniquely by a group homomorphism, so that it
takes prescribed values on the $e_i$.

\VerdierSemanticAnchor{ch1-1-6-3}
\VerdierSemanticAnchor{ch1-1-6-3-definition}
\noindent\textbf{Definition 1.6.3.}\quad
Let $e_1,\ldots,e_n$ be a basis of $\mathbb Z^n$. The
\emph{Koszul class relative to the basis $e_1,\ldots,e_n$} is the
unique element of $H^2(\mathbb Z^n,\mathbb Z/2\mathbb Z)$ whose
restrictions to the groups
$H^2(\mathbb Z^n_{i,j},\mathbb Z/2\mathbb Z)$ are nonzero.

\VerdierSemanticAnchor{ch1-1-6-3-basis-dependence}
Note that the definition of the Koszul class involves the choice of a
basis of $\mathbb Z^n$. The Koszul classes associated with two distinct
bases are distinct in general.

\VerdierSemanticAnchor{ch1-1-6-3-filtration}
Let $e_1,\ldots,e_n$ be a basis of $\mathbb Z^n$ ordered according to
the natural order on $[1,n]$. Let $\mathbb Z^n(\leq i)$ (respectively,
$\mathbb Z^n(\geq i)$) denote the subgroup of $\mathbb Z^n$ generated
by the $e_j$ with $j\leq i$ (respectively, $j\geq i$).
